% Generated draft: requires mathematical and source review.
\documentclass[11pt]{article}
\usepackage[margin=25mm]{geometry}
\usepackage{fontspec}
\setmainfont{Latin Modern Roman}
\usepackage{amsmath,amssymb,mathtools,unicode-math}
\setmathfont{Latin Modern Math}
\usepackage{enumitem,needspace,xurl,hyperref}
\setlength{\parindent}{0pt}
\setlength{\parskip}{5pt}
\setlength{\emergencystretch}{3em}
\newcommand{\N}{\mathbb{N}}
\newcommand{\Z}{\mathbb{Z}}
\newcommand{\Q}{\mathbb{Q}}
\newcommand{\R}{\mathbb{R}}
\newcommand{\C}{\mathbb{C}}
\newcommand{\F}{\mathbb{F}}
\newcommand{\A}{\mathbb{A}}
\newcommand{\PP}{\mathbb{P}}
\newcommand{\E}{\mathbb{E}}
\newcommand{\eps}{\varepsilon}
\newcommand{\dd}{\,\mathrm{d}}
\begin{document}
\section*{UnsolvedMath ID: 800232}
\textbf{Source problem code:} catalogue entry
\subsection{Definitions and mathematical statement}
Let \(M\) be a compact real analytic CR manifold with \(\dim_{\R}M\geq 5\), assumed pseudoconvex as in the source. Its CR structure is a real analytic complex subbundle \(T^{0,1}M\subset \C TM\) satisfying
\[
T^{0,1}M\cap\overline{T^{0,1}M}=\{0\},
\qquad
[\Gamma(T^{0,1}M),\Gamma(T^{0,1}M)]\subseteq\Gamma(T^{0,1}M).
\]
Write \(\Lambda^{0,q}M=\Lambda^q(T^{0,1}M)^*\), and let \(\bar\partial_b\) denote the tangential Cauchy--Riemann operator on these bundles. For auxiliary Hermitian inner products and a positive density defining the formal adjoint \(\bar\partial_b^*\), the Kohn Laplacian on \((0,1)\)-forms is
\[
\Box_b=\bar\partial_b\bar\partial_b^*+\bar\partial_b^*\bar\partial_b.
\]
Here \(\mathcal D'(M;\Lambda^{0,1}M)\) denotes distributional \((0,1)\)-forms, and \(C^\omega(M;\Lambda^{0,1}M)\) denotes real analytic \((0,1)\)-forms.

The source asks whether \(\Box_b\) is globally real analytic hypoelliptic on \(M\), suggesting \((0,1)\)-forms as the degree to consider. In that degree, the question is whether every such \(M\) satisfies
\[
\forall u\in\mathcal D'(M;\Lambda^{0,1}M),\qquad
\Box_bu\in C^\omega(M;\Lambda^{0,1}M)
\ \Longrightarrow\ 
u\in C^\omega(M;\Lambda^{0,1}M).
\]
This implication concerns analyticity on all of \(M\); no orthogonality to \(\ker\Box_b\) is imposed in the supplied question.

The source does not explicitly specify the CR codimension, the precise pseudoconvexity convention, or the auxiliary inner products and density used to define \(\Box_b\). In the usual hypersurface-type interpretation, \(\dim_{\R}M=2n+1\) with CR dimension \(n\geq2\), and pseudoconvexity means semidefiniteness of the Levi form with the appropriate local choice of sign. That interpretation is conditional and is not an additional supplied hypothesis. The intended geometric class and the regularity and quantification of the auxiliary choices require clarification before the question determines a fully specified operator problem.
\subsection{Short English statement}
For a compact real analytic pseudoconvex CR manifold of real dimension at least five, must every distributional \((0,1)\)-form whose Kohn Laplacian is real analytic on the whole manifold itself be real analytic on the whole manifold? The source leaves the CR codimension, pseudoconvexity convention, and auxiliary choices defining the operator implicit.
\subsection{Sources}
\begin{itemize}
\item Local UnsolvedMath registry: permanent ID \(800232\), exact source code \texttt{catalogue entry}, title ``Several Complex Variables Problem List --- Problem 29.'' The supplied record labels the problem open; this is a registry label, not an independently verified status assessment.
\item Buzzard, Fornaess, and Krantz, \emph{The Several Complex Variables Problem List} (1995), Problem 29. Supplied source URL: \url{https://arxiv.org/abs/math/9509202}. The record attributes this individual question to D. S. Tartakoff and supplies the historical contact \texttt{U22393@UICVM.bitnet}.
\item Supplied provenance: extraction method \texttt{arxiv\_latex}; source access date 2026-10-16; upstream revision \texttt{d371ee917ac9647fbb06cd6fc529deec6f3dda82}. These are supplied provenance details; no independent source access or literature verification is asserted.
\end{itemize}
\end{document}
